\documentclass[10pt,twocolumn]{article}

\usepackage[margin=0.75in]{geometry}
\usepackage{amsmath}
\usepackage{booktabs}
\usepackage[expansion=false]{microtype}
\usepackage{url}
\usepackage[hidelinks]{hyperref}
\usepackage[numbers,sort&compress]{natbib}

\hypersetup{
  pdftitle={Governance as Inference-Time Capability Substitution},
  pdfauthor={Layer 1 Labs},
  pdfsubject={GovernanceBench model-substitution preprint},
  pdfkeywords={software engineering agents, LLM governance, model routing,
    token efficiency, reproducibility}
}

\title{Governance as Inference-Time Capability Substitution:\\
Reducing Tokens per Correct Software-Engineering Task with Deterministic
Requirements and Verification}
\author{Layer 1 Labs\\
\url{https://github.com/layer1labs/specsmith}}
\date{Preprint --- July 2026}

\newcommand{\specsmith}{\textsc{Specsmith}}
\newcommand{\tpca}{\mathrm{TPCA}}
\newcommand{\cop}{\mathrm{CoP}}

\begin{document}
\maketitle

\begin{abstract}
Language-model comparisons usually hold the agent scaffold fixed and vary the
model. We study the complementary systems question: can deterministic
governance compensate for measured model capability in repository-level
software tasks? We introduce GovernanceBench, an executable benchmark and
matched $2\times2$ design in which a lower-tier and frontier model each run
with an ungoverned agent loop and with \specsmith{} FULL governance. The
controller exposes a stable five-tool interface, retrieves only
requirement-linked context, enforces visible tests without exposing hidden
acceptance oracles, and audits every trace. Our primary efficiency measure is
tokens per correct answer; correctness is a hard non-inferiority constraint.
In the preregistered eight-task, ten-repetition experiment, GPT-5.6 Terra with
FULL governance passed 80/80 cells at 11,737 tokens per correct answer, while
GPT-5.6 Sol in the raw loop passed 65/80 at 28,023. The TPCA ratio was 0.419
(fixed-suite 95\% interval 0.358--0.485; task-cluster interval 0.212--0.708).
Both preregistered mixed-suite substitution gates passed. In a prespecified
coding-only sensitivity, Terra FULL passed 60/60 at 15,649 TPCA versus Sol raw
55/60 at 30,250, but its correctness lower bound was $-5.7$ percentage points,
just outside the $-5$-point non-inferiority margin. Published 8B and mini-model
routes failed admission. In a separately frozen real-repository lane, three
complete 27B--32B routes also failed correctness and three pinned 20B--35B
routes were censored by provider incompatibility or unavailability; no route
earned repeated promotion. On corrected pinned-upstream T30, Sol FULL passed
8/10 at 105,523 TPCA versus raw Sol at 5/10 and 234,526, while Terra FULL
passed 2/10 at 362,650. The within-model joint gates remained unconfirmed, and
the lower-tier substitution result did not transfer. The evidence supports
task-conditional substitution for the exact API routes and benchmark
distribution, not universal model equivalence or a claim that small models
replace frontier models.
\end{abstract}

\section{Introduction}

Repository-level coding agents must localize requirements, edit multiple
files, run tools, interpret failures, and decide when evidence is sufficient.
Benchmarks such as SWE-bench expose this interaction complexity
\citep{jimenez2024swebench}, while SWE-agent shows that the agent-computer
interface itself can materially change performance \citep{yang2024sweagent}.
Consequently, ``which model is best?'' is incomplete: practical outcomes depend
on the joint system of model, context policy, tools, validators, and stopping
rules.

This paper asks whether a governance controller can act as
\emph{inference-time capability substitution}. The target is not to make a
lower-tier route intrinsically equivalent to a frontier model. It is to determine
whether a lower-tier model embedded in a deterministic, evidence-seeking
controller can preserve task correctness while consuming fewer tokens or
dollars than a frontier model in an ungoverned loop.

The question connects three literatures. Interactive reasoning methods ground
generation in external actions and feedback \citep{yao2023react,
shinn2023reflexion,madaan2023selfrefine}. Simpler software-agent pipelines can
outperform more elaborate loops \citep{xia2024agentless}. Model cascades and
routers exploit heterogeneous cost--quality tradeoffs \citep{chen2023frugalgpt,
ong2024routellm}. \specsmith{} differs by routing not only between models but
between probabilistic generation and deterministic governance operations:
requirement lookup, scoped retrieval, test execution, safety refusal, evidence
compression, and stopping.

We make four contributions:
\begin{enumerate}
  \item a matched $2\times2$ experimental design separating model capability
        from governance lift;
  \item a primary metric, tokens per correct answer, coupled to a preregistered
        correctness non-inferiority gate;
  \item deterministic fixed-suite and task-cluster uncertainty analysis plus
        machine-readable claim gates; and
  \item a versioned compact evidence artifact containing per-cell outcomes,
        workflow provenance, negative results, and a post-run weakness audit,
        backed by versioned tasks, validators, and hidden acceptance oracles.
\end{enumerate}

\section{Related Work}

\paragraph{Executable software-engineering benchmarks.}
SWE-bench evaluates issue resolution against executable repository tests
\citep{jimenez2024swebench}. Static benchmark contamination and memorization
are material threats \citep{liang2025illusion}; SWE-bench-Live addresses
freshness and repository diversity with continuously updated tasks and
containerized environments \citep{zhang2025swebenchlive}. GovernanceBench is
smaller and synthetic by design. It cannot estimate population-level coding
ability, but its fixtures make controller ablation, exact counterfactual
matching, and trace-level token accounting practical.

\paragraph{Agent interfaces and feedback.}
SWE-agent attributes substantial gains to a model-oriented interface
\citep{yang2024sweagent}. ReAct interleaves reasoning with environmental
actions \citep{yao2023react}; Reflexion and Self-Refine use iterative feedback
without parameter updates \citep{shinn2023reflexion,madaan2023selfrefine}.
Agentless demonstrates that constrained localization, repair, and validation
can be competitive at low cost \citep{xia2024agentless}. These findings
motivate a compact tool surface and controller-owned validation rather than a
large catalogue of agent skills.

\paragraph{Cost-aware model selection.}
FrugalGPT and RouteLLM demonstrate that heterogeneous models can be cascaded or
routed to improve cost--quality tradeoffs \citep{chen2023frugalgpt,
ong2024routellm}. Our unit of analysis is a complete software task, and our
router signal is not a learned preference score. Admission and post-run audit
use executable correctness, requirements, scope boundaries, token receipts,
and task-specific historical envelopes.

\paragraph{Reproducibility and variance.}
Benchmark point estimates can be dominated by uncontrolled variation
\citep{bouthillier2021variance}. Recent analysis of 640 LLM-for-software-
engineering papers reports persistent artifact, environment, and versioning
gaps \citep{siddiq2025reproducibility}. We therefore repeat every release cell,
pin task and controller revisions to one commit, reject missing provider cells,
record model routes and tool-schema hashes, and publish interval estimates
rather than relying on means alone.

\section{System}

\subsection{Governance controller}

\specsmith{} FULL wraps a language model with six mechanisms:
\begin{enumerate}
  \item \textbf{Requirement contracts}: accepted requirements, linked tests,
        scope, risk, and epistemic constraints are loaded before edits.
  \item \textbf{Bounded context}: versioned requirement-linked files may be
        preloaded; all other content is retrieved just in time.
  \item \textbf{Stable interface}: every model turn sees the same five tools:
        batched and scalar read/write operations plus completion.
  \item \textbf{Deterministic gates}: ambiguous or destructive tasks can stop
        without an LLM call; coding completion requires current public
        validation evidence.
  \item \textbf{Evidence compression}: unchanged reads become digest receipts,
        superseded observations are removed, and only active milestone state
        remains in long-horizon context.
  \item \textbf{Hidden scoring}: acceptance oracles run only after the model
        stops and never supply repair hints.
\end{enumerate}

The design intentionally excludes broad repository preloads, mutable tool
schemas, evaluator leakage, and unbounded self-reflection. A deterministic
post-run audit identifies correctness, scope, retrieval, repair, tool-schema,
and cost weaknesses and selects the next admissible experiment.

\subsection{Promotion funnel}

New model routes receive one long-horizon governed admission. Controller
changes receive one cell on each affected path. Successful candidates advance
to a four-task, five-repetition $2\times2$ screen. Only then may the complete
eight-task, ten-repetition substitution profile run. This funnel prevents a
failed candidate from becoming persuasive merely through additional spend.

\section{GovernanceBench}

\subsection{Tasks and conditions}

The release grid contains eight tasks: two feature/bug tasks, two deterministic
clarification or safety gates, three cross-file feature/refactor tasks, and one
long-horizon polyglot incident console spanning Python, Go, TypeScript/React,
JSON Schema, browser tests, and documentation. Six tasks require code changes;
two assess whether the system declines or clarifies unsafe or underspecified
work.

The substitution experiment crosses two models with two conditions:
\begin{center}
\begin{tabular}{ll}
\toprule
Cell & Purpose \\
\midrule
Lower-tier raw & lower-tier capability baseline \\
Lower-tier FULL & candidate governed system \\
Frontier raw & substitution target \\
Frontier FULL & frontier governance control \\
\bottomrule
\end{tabular}
\end{center}

All four cells use identical task fixtures, repetition identifiers, provider
APIs, temperature policy, turn cap, completion cap, and hidden oracles.
Provider errors invalidate a cell rather than counting as free failures.

\subsection{Metrics}

For runs $i=1,\ldots,n$, with token count $t_i$, cost $c_i$, and correctness
$y_i\in\{0,1\}$, we report
\begin{align}
\tpca &= \frac{\sum_i t_i}{\sum_i y_i}, &
\cop &= \frac{\sum_i c_i}{\sum_i y_i}.
\end{align}
This formulation charges failed attempts to the expected cost of obtaining a
correct answer. Raw tokens are primary because provider prices and cache
discounts change; estimated dollars are secondary.

Correctness is evaluated by deterministic lint, public tests, task-specific
contract validators, and an evaluator-isolated acceptance oracle. We also
record turns, wall time, files read and written, scope expansion, first-pass
rate, provider-reported cache tokens, and a stable tool-schema hash.

\subsection{Hypotheses and claim gates}

The release profile and thresholds were committed at revision
\texttt{75a8c79} before workflow dispatch. The primary comparison is
lower-tier FULL minus frontier raw:
\begin{description}
  \item[H1: correctness non-inferiority.] The lower 95\% bound on the pass-rate
        difference is at least $-5$ percentage points.
  \item[H2: token superiority.] The upper 95\% bound on the TPCA ratio is below
        1.
  \item[H3: cross-task robustness.] H1 and H2 also hold when task identifiers,
        not only repetitions, are resampled.
  \item[H4: governance lift.] Lower-tier FULL improves correctness or TPCA
        relative to the same lower-tier raw model.
\end{description}

We use Wilson score bounds for binomial rates \citep{wilson1927probable} and a
deterministically seeded, 10,000-sample hierarchical bootstrap
\citep{efron1979bootstrap}. The fixed-suite analysis resamples repetitions
within each task. The task-cluster analysis additionally resamples task IDs;
it tests robustness to resampling the benchmark tasks but cannot establish
performance on fresh repositories. A result can therefore support fixed-suite
substitution while failing cross-task substitution. No result supports
universal model equivalence.

\section{Results}

% Generated by scripts/govern_bench/render_preprint_results.py; do not edit.
\begin{table*}[t]
\centering
\caption{Preregistered eight-task, four-cell, n=10 substitution replication. Workflow 30210886840, commit \texttt{75a8c7911187}. TPCA is tokens per correct answer and CoP is estimated API cost per passing task.}
\label{tab:release-results}
\begin{tabular}{lrrr}
\toprule
System & Correct & TPCA & CoP (USD) \\
\midrule
gpt-5.6-terra raw & 56/80 & 45,915 & 0.1642 \\
gpt-5.6-terra + Specsmith FULL & 80/80 & 11,737 & 0.0602 \\
gpt-5.6-sol raw & 65/80 & 28,023 & 0.1511 \\
gpt-5.6-sol + Specsmith FULL & 80/80 & 8,082 & 0.0752 \\
\bottomrule
\end{tabular}
\end{table*}

\paragraph{Preregistered release inference.}
The lower-tier FULL/frontier-raw TPCA point ratio was 0.419. Its fixed-suite 95\% interval was 0.358--0.485; the task-cluster interval was 0.212--0.708. The corresponding correctness-difference intervals were 3.3--34.1\,pp and -4.4--54.7\,pp. The machine-readable release gate passed. Fixed-suite substitution was supported, and cross-task substitution was supported.

\paragraph{Coding-only sensitivity.}
After excluding deterministic clarification and safety gates, lower-tier FULL passed 60/60 coding cells at 15,649 TPCA; frontier raw passed 55/60 at 30,250 TPCA. The TPCA ratio was 0.517, with fixed-suite 95\% interval 0.437--0.608 and task-cluster interval 0.328--0.870. The correctness-difference intervals were -5.7--26.1\,pp and -6.0--31.8\,pp. Coding-only fixed-suite substitution was not supported; this sensitivity analysis was specified for interpretation after the mixed-suite release gate and is not relabelled as the primary endpoint.


The primary mixed-suite point estimate corresponds to 58.1\% lower TPCA and
60.1\% lower estimated cost per correct answer for Terra FULL than Sol raw.
The task-cluster correctness lower bound was $-4.4$ percentage points, clearing
the preregistered $-5$-point margin by 0.6 points. The coding-only sensitivity
went in the same observed direction but did not clear that margin; deterministic
clarification and safety gates therefore matter to the strength of the primary
claim. We report both endpoints rather than relabelling the sensitivity after
observing the result.

The four-task screen demonstrates why both uncertainty views are needed. Its
point TPCA ratio is 0.653. The fixed-suite 95\% interval is 0.539--0.784, while
the task-cluster interval is 0.339--1.396. Terra FULL is less token-efficient
than Sol raw on T1 and T10 but substantially more efficient on T13 and T28.
The aggregate is therefore useful routing evidence, not a universal
replacement result.

\section{Trace-Guided Ablations}

The benchmark audit converts failures and expensive correct runs into bounded
controller experiments. Three examples illustrate the loop:
\begin{itemize}
  \item A T10 controller denied reads of an imported model and dependency
        manifest. Adding exactly those files reduced one admission from 63,128
        tokens and 11 turns to 9,305 tokens and two turns.
  \item Repeated T11 first actions justified a bounded preload, reducing a
        diagnostic from 23,472 tokens and six turns to 11,210 tokens and two
        turns.
  \item Applying the same idea to T13 increased tokens from a 14,424-token
        broad-run mean to 16,190 in admission. The preload was removed.
\end{itemize}
The negative T13 result is important: more context is not monotonically
beneficial, and governance optimization itself requires admission controls.

Open-model admissions provide a second negative control. Qwen3.6-35B-A3B
completed the long-horizon task with a stable tool schema but used 67,701
tokens, 3.87 times the governed frontier anchor. It was not promoted to repeated
spend. Correctness alone is insufficient for an efficiency claim.

A published-size Llama 3.1 8B route failed three governed T28 admissions.
Strict recovery of its text-serialized function calls removed a serving
protocol ambiguity, but the final route still issued 20 actions serially,
exhausted the turn cap, and failed at 98,543 tokens. GPT-4o mini also failed
T28 at 107,896 tokens after focused-boundary repair. These admissions block
the stronger marketing claim that an 8B or mini model replaces a frontier
model. They do not negate the separately tested lower-tier Terra system, whose
parameter count and architecture are not inferred.

A later reasoning-capable cohort covered exact managed routes from 20B to 32B.
GPT-OSS-20B, Qwen3-Coder-30B-A3B, GLM-4.7-Flash, and Qwen3-32B failed their
governed T28 admissions. Qwen3.6-27B initially stopped after narrating its next
milestone; a bounded controller repair then produced a correct public-and-hidden
cell at 72,255 tokens and 14 model turns. This establishes n=1 capability for
the published 27B route, but not efficient substitution: the cell used 4.13
times the 17,501.7-token governed Sol envelope and was not promoted to repeated
spend. The combined evidence separates ``can finish once'' from ``replaces a
frontier system efficiently.''

A subsequent exact-boundary confirmation also passed at 77,776 tokens and 14
turns. It is not pooled with the first cell because the controller commit
changed between runs. Its higher 4.44-times frontier-envelope ratio shows that
27B capability can recur while preserving the negative efficiency conclusion.

A subsequent controller-protocol isolation produced a material within-route
improvement. Automatic scalar-parallel file calls completed the same T28 public
and hidden gates in six turns and 26,850 tokens, 62.8 percent below the original
correct admission. The cell is still an n=1 diagnostic and consumes 1.53 times
the 17,501.7-token governed Sol envelope. Required-tool, write-only,
context-compaction, and independent-validator-authority variants did not beat
it. The result therefore supports editing-interface optimization, not a
frontier-replacement claim.

A later native-interface series isolated provider-visible editing structure.
Managed Novita routes for Qwen3-Coder-Next and Qwen3-Coder-480B failed the same
task. A fixed-scalar milestone bundle restored Qwen3.6-27B correctness at 71,090
tokens and 13 turns, but consumed 4.06 times the governed Sol envelope and 2.65
times the scalar-parallel winner. Exact single-hunk edits remained near 27,000
tokens but could not express a coupled import-and-annotation repair atomically.
The resulting bounded three-hunk patch interface passed local safety and routing
tests; its live runtime was censored by a provider 504, and the single retry
failed route availability before execution. Both censored attempts are excluded
from correctness and efficiency inference. These results show that tool-serving
structure can alter correctness while leaving the broad frontier-replacement
claim unsupported.

A follow-up used the same atomic interface with the hosted
Qwen3-Coder-30B-A3B route after an exact structured-tool probe. The baseline
failed at 123,384 tokens, 20 turns, and four of ten files. A scoped variant
removed model-owned rereads and reduced the failed attempt to 34,892 tokens,
10 turns, and 29.86 seconds---reductions of 71.7 percent in tokens, 73.9
percent in estimated cost, and 78.3 percent in wall time. It still failed the
independent oracle, so this is a failure-cost result rather than a TPCA or
substitution result. Requiring a tool on every turn regressed to 81,648 tokens
and repeated one three-action batch across eleven turns. We reject required
tool choice and add a content-free batch-repeat guard that would stop the
observed sequence after its fourth identical batch.

We separately preregistered an ephemeral endpoint for the official FP8
checkpoint, vLLM 0.24.0, and the native \texttt{qwen3\_xml} parser
\citep{qwen3coder30b,vllmtoolcalling,hfendpoints}. The first attempt was
censored before provisioning because the credential lacked endpoint-write
permission. After correcting that permission, two endpoint runs passed an exact
required-tool probe and confirmed pause and deletion. Atomic and scoped
automatic-tool cells still failed at 34,146 and 53,451 tokens. A scoped
required-tool cell failed at 181,884 tokens and 20 turns after changing the
failure mode from premature narration to completion truncation and milestone
fragmentation. The two endpoint receipts report approximately \$0.430977
combined infrastructure cost. Thus native parser compatibility removes a
serving-stack ambiguity but is insufficient for correctness, and globally
required tool choice is rejected for this model and task.

We then tested the newer Qwen3.6-35B-A3B-FP8 checkpoint on one A100 with vLLM
0.24.0, a 65,536-token context, the native \texttt{qwen3\_coder} tool parser,
and the \texttt{qwen3} reasoning parser. Both cells passed an exact tool probe
and recorded successful endpoint pause and deletion. Generic required-tool
recovery failed T29 after one of four milestones, 30,287 tokens, a length stop,
and two empty continuations. Exact named-tool recovery advanced through three
milestones but failed at 105,087 tokens and the 20-turn cap. Deterministic
sanitation removed 96 duplicate actions, and bounded malformed-patch recovery
still left Go dependency and TypeScript query repairs incomplete. The latter
cell consumed 5.97 times the correct Terra T29 v7 token anchor without a correct
answer. This valid native-route failure rejects repetition and shows that parser
compatibility plus governance recovery did not substitute for model capability
on the fresh polyglot task.

The same controls were then exercised through the native Responses tool route
of the stronger GPT-5.6 Sol system. Replacing eight fixed nullable file slots
with one bounded strict array of file objects preserved all-or-nothing
milestone writes while reducing provider-visible schema overhead. An
independent $n=10$ T28 confirmation passed public and hidden validation in all
ten rows, completed in five model turns with 100 percent first-pass success,
and averaged 17,907 tokens per correct answer (1.77 percent coefficient of
variation). This was 7.2 percent below the preceding stable-schema native
controller, but 2.3 percent above the 17,502-token release anchor. The result
supports native structured editing as a material controller optimization on
this task and route; it does not establish a universal model or task-family
advantage.

We then froze v4 and changed only the model route to GPT-5.6 Terra. After a
visible-validator repair, T28 passed five of five screening rows, but its
independent ten-row confirmation passed only nine. The failed row repeatedly
targeted the first milestone until the bounded loop guard stopped execution,
leaving three product boundaries incomplete. The sample consumed 176,832
tokens and yielded 19,648 tokens per correct answer, so the audit rejected
promotion. We next introduced T29, an independent synthetic release-control
repository with a different domain, seven-field contract, starter,
validators, and hidden oracle while preserving the four milestone sizes and
frozen controller. T29 passed five of five rows at 22,038 mean tokens per
correct answer and \$0.1299 mean estimated cost, but every row entered App
empty-state repair and one invoked loop recovery. Re-audit showed the original
CSS-hotspot label was a parser artifact from a relative import in supplied App
content. The public validator required the literal word ``empty'' even when
the hidden oracle accepted the structural zero-length/no-release state.

A trace-derived v5 package aligned that check, added CSS-only milestone-three
validation for focus, disabled actions, and responsive layout, and bounded
unchanged repair streaks at two. Admission workflow 30453721376 passed
first-pass at 17,407 tokens. The earned five-row workflow 30454018932 passed
five of five at 18,055 mean tokens per
correct answer, \$0.1199 mean estimated cost, 5.2 turns, 80\% first-pass, and
zero loop recoveries. One row repaired a missing Playwright visibility
assertion in one patch. Relative to v4, mean tokens fell 18.1\%, cost 7.7\%,
and turns 18.8\%; because the task/validator contract and recovery controller
changed together, this is a package-level delta rather than a single-factor
causal estimate. The admitted independent ten-row workflow 30457360342 passed
all ten rows at 20,099.9 mean tokens per correct answer, \$0.127686 mean
estimated cost, 5.7 turns, and 9.8\% coefficient of variation. First-pass
completion was only 30\%: six rows repaired a missing Playwright visibility
assertion and one repaired an invalid model-authored pytest assertion. Relative
to the favorable five-row estimate, mean tokens rose 11.3\%, cost 6.5\%, and
turns 9.6\%. The audit therefore selects \texttt{optimize\_and\_rerun}, not
clean promotion. This supplies exact-route, synthetic release-sized evidence
while preserving the negative conclusion on real-repository generalization.

Two further trace-derived packages tested whether those repair receipts could
be moved into milestone-boundary constraints. Versioned v6 made the existing
Playwright visibility criterion explicit and removed the recurring UI repair,
but two of five rows still repaired milestone-one Python toolchain mistakes.
Versioned v7 additionally made those observed constraints explicit. Its
admission and five-row screen passed every row first-pass. The independent
ten-row workflow 30547516220 then passed public and hidden validation in all
ten rows first-pass at 17,589.9 mean tokens per correct answer, \$0.120940 mean
estimated cost, five turns, and 1.62\% coefficient of variation. Relative to
the v5 ten-row sample, mean tokens fell 12.5\%, estimated cost 5.3\%, and turns
12.3\%, while first-pass completion rose 70 percentage points. This establishes
a clean repeated result for the exact synthetic T29, Terra Responses, and FULL
controller package; it still does not establish transfer to a fresh real
repository.

We therefore froze protocol
\texttt{GB-PREPRINT-2026-07-30-V1} before testing a byte-hashed snapshot of
Pallets ItsDangerous at commit
\texttt{672971d66a2ef9f85151e53283113f33d642dabd}. The untouched upstream suite
passed 297 tests; the starter failed both a public key-rotation validator and
an evaluator-isolated oracle. Six exact hosted 20B--35B routes then received
one T30 FULL admission with a 30,000-token correctness gate. Qwen3.6-27B,
Qwen3-32B, and Qwen3-Coder-30B produced complete failures at 52,465, 148,710,
and 28,821 tokens. The first stopped after one of three milestones, the second
exhausted twelve turns after writing all declared boundaries, and the third
returned an empty post-tool continuation before completing a milestone.
GPT-OSS-20B rejected the frozen named tool choice after 8,612 partial tokens;
GLM-4.7-Flash produced a 504 and then a bounded timeout after 7,873 partial
tokens; Qwen3.6-35B timed out in two live probes before a scored cell. Complete
failed cells consumed 229,996 tokens and \$0.080464, while censored calls add
16,485 observed tokens and \$0.002487. No complete route passed the independent
public project-test gate, so the preregistered sequential rule prohibited n=5
or n=10 promotion. This is a negative real-repository admission result, not
evidence that governance alone makes a mid-sized route frontier-equivalent.

The matched V1 workflow then exposed an evaluator defect before publication:
its T30 hidden oracle required a prose license label absent from the untouched
pinned file. T30 was also misclassified as standard horizon, so FULL could
request completion before every milestone file existed. We invalidate all 60
V1 T30 rows and retain their expenditure only as invalidated experimental
spend. The open-model rejection remains valid because each complete route
independently failed the public project-test gate. V1 T28 and T29 use separate
fixtures and oracles and remain protocol-labeled evidence, although two
Terra-raw T28 calls ended in bounded provider timeouts.

Before further paid calls, we froze
\texttt{GB-PREPRINT-2026-07-30-V2}. It reruns only T30 at n=10 for Terra and
Sol under raw, Cursor-style, and FULL conditions, verifies the license by its
pinned digest, enforces long-horizon milestone completion, and disables
provider retries. This recovery separates evaluator correction from controller
optimization and prevents V1 outcomes from entering the corrected inference.

V2 workflow 30589098641 completed all 60 cells without a provider retry,
error, skip, or censored observation. Sol FULL passed 8/10 at 105,523 TPCA and
\$0.52046 per pass, compared with raw Sol at 5/10, 234,526, and \$0.66978 and
Cursor-style Sol at 1/10, 1,027,057, and \$2.94382. Failed rows consumed
26.6\% of Sol FULL tokens, versus 47.7\% raw and 85.2\% Cursor-style. Terra
FULL passed 2/10 at 362,650 TPCA, compared with 1/10 for both Terra controls.
The Sol FULL-versus-raw TPCA interval favored FULL at 0.222--0.656, but its
correctness-difference interval was $-45.1$ to $+87.4$ percentage points; all
four within-model joint gates remained false.
Separately, Terra FULL failed the prespecified cross-model substitution
comparison against raw Sol, which passed 5/10 at 234,526 TPCA. Some paired
bootstrap resamples contained no correct
Sol-raw answer, making the upper TPCA-ratio bound non-estimable; we serialize
that bound as null and fail both substitution gates closed.

Trace audit localized the residual failures. Sol FULL's two failures exhausted
the turn cap after incomplete timed-value and documentation work. Terra FULL
spent 76.9\% of its tokens on failures: two cells stopped on prose after one
turn, others serialized actions or exhausted the cap, and two satisfied the
hidden rotation contract while regressing the unchanged upstream suite. The
dual public/hidden gate prevented those partial solutions from receiving
credit.

\section{Threats to Validity}

\paragraph{Construct validity.}
GovernanceBench measures repository-task execution under one controller, not
general intelligence, developer productivity, code maintainability, or factual
truth. Synthetic fixtures improve experimental control but may reward explicit
contracts and validators more than naturally occurring repositories.

\paragraph{Internal validity.}
Proprietary model weights, serving stacks, and sampling implementations are not
controlled. Repetition IDs do not guarantee identical provider randomness.
Dollar costs are versioned estimates, not invoices. We mitigate these threats
with same-commit matched cells, exact model-route labels, raw token receipts,
fail-closed provider handling, and paired uncertainty.

\paragraph{External validity.}
Eight tasks are too few for an unconditional market-wide replacement claim.
Task-cluster intervals expose this limitation but cannot manufacture diversity.
The benchmark is public and may eventually be contaminated. T30 adds one
pinned independent upstream repository, but broader replication on fresh,
real repositories such as SWE-bench-Live remains necessary
\citep{zhang2025swebenchlive}.

\paragraph{Researcher degrees of freedom.}
Earlier runs informed controller development, so the study includes adaptive
optimization. The decisive release profile, thresholds, commit, and selected
model pair were fixed before its run. We publish negative admissions and do
not pool incompatible commits.

\section{Operational Implications}

The evidence motivates a guarded router rather than a single default model.
Deterministic gates should handle tasks that do not require generation.
Lower-tier governed models should receive bounded, validator-rich work for
which they have passed admission. Frontier governed models remain appropriate
for ambiguous localization, novel architectures, and task families whose
lower-tier confidence bound crosses the correctness or efficiency threshold.
Routing is therefore based on verified marginal value, consistent with the
cost-aware motivation of FrugalGPT and RouteLLM
\citep{chen2023frugalgpt,ong2024routellm}, but grounded in executable software
evidence.

\section{Reproducibility and Ethics}

The repository publishes the harness, fixtures, controller, exact workflow
inputs, analysis code, compact per-cell measurements, and SHA-256 evidence
manifest. Full traces are retained with the cited GitHub Actions runs according
to the repository's configured artifact-retention period; the compact evidence
and source digests are versioned permanently. The paper reports model API
identifiers but does not infer unpublished parameter counts or architecture.
Hidden acceptance tests are isolated from model feedback. Safety tasks reward
clarification or refusal rather than destructive execution. No human-subject
data are collected.

The corrected T30 bundle contains the 60 canonical rows, both source-artifact
digests, workflow and commit identities, and V2 protocol SHA-256. A separate
machine-readable record invalidates only V1 T30 while preserving independently
valid T28/T29 strata and the two censored Terra T28 cells. The preprint has not
been submitted to an external server.

\section{Conclusion}

Governance can move work from probabilistic model reasoning into deterministic
requirements, retrieval, verification, and stopping. On the evaluated mixed
task distribution, this substantially reduced the tokens required for correct
software changes and allowed the governed Terra route to clear the
preregistered substitution gates against raw Sol. The coding-only gate remained
inconclusive by 0.7 percentage points, and the published-size admissions
failed. On the independent T30 repository, FULL produced favorable Sol point
estimates but did not clear the joint within-model uncertainty gate, and Terra
did not substitute for raw Sol. The scientifically defensible result is
task-conditional system substitution with explicit confidence and routing
gates---not that smaller models universally replace frontier models.

\bibliographystyle{plainnat}
\bibliography{references}

\end{document}
